\section[JVM network contracts]{JVM network and simulation components: making exposure and intervention clocks explicit}
\label{sec:exec25-java-network}

A bank constructing borrower-network experiments on the JVM must separate the graph representation from the causal experiment. The graph can encode households, suppliers, branches or common employers; its edges may mean influence, contractual exposure or merely similarity. A simulator can generate outcomes after an intervention has been specified, but its output does not identify that intervention from observed bank data. This comparison evaluates JGraphT, GraphStream, MASON and Repast Simphony as components for graph construction and public experiments. Admission as a component is separate from qualification of a causal estimator.~\citep{exec25_jnet_jgrapht,exec25_jnet_graphstream,exec25_jnet_mason,exec25_jnet_repast}

The inspected revisions precede the 7 October 2026 client cutoff. JGraphT's selected tree declares 1.6.0-SNAPSHOT and Java 21; the release snippet in its README is not the selected build. GraphStream declares 2.1.0-SNAPSHOT in the selected November 2021 tree. The MASON review binds the July 2026 source revision, while the Repast core manifest identifies 2.12.0 at a September 2026 revision. Fifty technical, license and configuration originals plus two candidate-screening originals were saved and authenticated. Eight bounded repository queries return 71 distinct leads. Some query results have uninspected tails, and the returned hits include computing, radio, genetics and distributed-system meanings of interference or causality. These results cannot establish absence of Java network-causal software.~\citep{exec25_jnet_choice}

One seemingly relevant partial-interference repository describes colocated server workloads, resource metrics and an SLA scheduler in its original README. It is not selected as evidence of a borrower-spillover estimator. A solar-innovation search lead has an empty root README and remains unresolved without technical inspection. The distinction prevents a search label from becoming a capability claim. No downloaded code, native JVM artifact, upstream test or simulator was executed; no theorem or benchmark is imported from papers mentioned in source comments.~\citep{exec25_jnet_choice}

\subsection{JGraphT: neighbor lists, graph distances and live subgraphs}

JGraphT provides distinct operations that a causal adapter should preserve. Its neighbor list traverses all incident edges and can include a vertex more than once in a multigraph; its neighbor set removes duplicate vertices. Predecessor and successor lists instead traverse incoming and outgoing edges. These are useful alternatives, not interchangeable encodings of a bank exposure. Whether repeated transactions count as repeated influence, summed exposure weight or a single relationship is a scientific decision.~\citep{exec25_jnet_jgrapht}

Consider the authored directed multigraph with two edges $a\to b$ and one edge $b\to c$, with treatment vector $(D_a,D_b,D_c)=(1,0,0)$. For node $b$, the mean treatment over the incoming list is one. Over the full incident-neighbor list it is $2/3$, because $a$ appears twice and $c$ once. Over the unique adjacent-vertex set it is $1/2$. A downstream model that sees only a field named neighbor exposure cannot distinguish these quantities. The proposed adapter should record direction, multiplicity, weighting, self-edge handling and normalization with the exposure operator.~\citep{exec25_jnet_choice}

The breadth-first shortest-path implementation traverses outgoing edges and increments distance by one. It therefore computes hops, regardless of stored edge weight. Its returned single-source object retains those hop distances; when reconstructing a path, the delegate sums the actual weights of selected edges. In an authored triangle with weights $w_{ab}=10$, $w_{ac}=w_{cb}=1$, the direct path has one hop and reconstructed weight ten, while the minimum-weight route has two edges and total weight two. The hop-distance query, the weight of the hop-selected path and the minimum-weight distance are three different quantities. An unreachable target has an infinite stored distance and no reconstructed path; it must not be silently converted to a zero-distance neighbor.~\citep{exec25_jnet_jgrapht,exec25_jnet_choice}

Subsetting also changes the question. In the path $a\to b\to c$, the full-graph distance from $a$ to $c$ is two. The induced graph on $\{a,c\}$ has no such path. Filtering the analysis population while retaining background network paths is therefore different from filtering the graph itself. This distinction matters when a bank observes outcomes only for a subset of borrowers but influence passes through other entities.~\citep{exec25_jnet_choice}

The selected subgraph class is a view over a backing graph. Its documentation and implementation distinguish listener-aware topology updates from a non-listenable base. Its edge-weight setter explicitly delegates to the backing graph. General wording that subgraph modifications do not affect the base must consequently be read as a topology statement, not a guarantee that all attributes are isolated. A scenario that changes edge weights in a subgraph can change the backing graph used by another scenario. The proposed contract uses immutable input identifiers and a deliberate deep-copy or versioned-view policy, with tests of both topology and attributes.~\citep{exec25_jnet_jgrapht}

JGraphT is a conditional direct-integration candidate for graph operators. The source supports explicit seeds or a supplied generator for its selected preferential-attachment generator, but repeatable generation still depends on vertex construction, iteration order, algorithm version and the graph's admissible edges. Synthetic topology is not authentic bank structure. The dual LGPL 2.1-or-later or EPL 2.0 declaration and complete local license texts support a concrete distribution review, while the selected snapshot, Java requirement and transitive components still need artifact-specific qualification. Porting should be considered only where the target runtime or numerical contract justifies the additional maintenance burden.~\citep{exec25_jnet_jgrapht,exec25_jnet_choice}

\subsection{GraphStream: ingestion policy changes the observed population}

GraphStream provides a different component: a graph that evolves through events. Its default edge-addition overload creates an undirected edge, while an explicit directed addition uses the specified source-to-target orientation. The DGS parser maps forward and reverse arrow tokens to corresponding directed events, and a step token becomes a numerical graph-step event. That step is not, by itself, the bank's event time, ingestion time or availability time. Those clocks require explicit fields and a documented conversion.~\citep{exec25_jnet_graphstream}

Error behavior is controlled by strict checking and automatic creation. A duplicate identifier may throw, return an existing compatible edge, or return null for incompatible endpoints. Missing endpoints may throw, produce null, or be created automatically. Edge rejection by a node similarly can throw or return null, with rollback of the first endpoint when necessary. These paths affect the observation contract: automatic creation can enlarge the population, while a discarded null can remove a relationship without telling the analyst. A causal ingestion adapter should return a receipt for each requested node or edge, including accepted, deduplicated, created, rejected and unresolved outcomes.~\citep{exec25_jnet_graphstream}

SingleGraph and MultiGraph differ at their node delegates. The selected SingleNode stores incoming and outgoing slots per opposite node and rejects conflicting additions; MultiNode retains a list of parallel edges. A read-only collection wrapper over that list does not freeze the underlying edge objects. A bank replay should therefore bind both the ordered event sequence and the resulting immutable graph snapshot, rather than retain a mutable collection reference as historical evidence. It should test repeated identifiers, reverse arrows, parallel obligations, unknown nodes and late-arriving attributes.~\citep{exec25_jnet_graphstream}

The proposed decision is conditional integration for dynamic ingestion and visualization, with causal exposure construction in a separately specified layer. The source's dual CeCILL-C or LGPL 3-or-later declaration requires a distribution choice and dependency review; it is not an unrestricted artifact clearance. No benefit follows from porting event handling unless the destination must preserve the same failure, ordering and time semantics. An agent choosing a permissive ingestion mode should not silently change the target population.~\citep{exec25_jnet_graphstream,exec25_jnet_choice}

\subsection{MASON: simulation state, update order and checkpoint identity}

MASON's network stores arbitrary node objects and edge information. Incoming and outgoing edge bags are internal mutable structures, and adjacency structures must be regenerated after topology changes. Its simple adjacency matrix is indexed from source to target and stores a single selected edge for a pair; a separate multigraph method is available when all parallel edges are needed. The selected weight getter returns a numeric or Valuable information object's value, and otherwise returns one. Its weight setter replaces the information object with a numeric value. A bank adapter should reject an unexplained unit-weight fallback where an edge actually represents a missing monetary amount or an uncertain relationship.~\citep{exec25_jnet_mason}

Copying the network does not automatically isolate a counterfactual scenario. The selected copy routine reuses node objects and passes the same edge-information references into new edge wrappers. A mutable customer state or edge-information object can therefore be shared across copies. The proposed scenario runner must define which state is copied, which is immutable, and which exogenous draws are intentionally coupled. It should test that changing a treatment or attribute in one scenario does not unintentionally mutate another. A full-state serialization checkpoint and a graph-wrapper copy solve different problems.~\citep{exec25_jnet_mason}

The scheduler gives equally consequential semantics. Lower integer orderings run first; equal-time, equal-order substeps are shuffled and then invoked sequentially. A newly scheduled event at the current time advances to the next representable time step. Thus a group of events with the same nominal timestamp is not necessarily a simultaneous update from the previous state. A simulator needs an explicit read-old/write-new protocol when that is the intended causal mechanism.~\citep{exec25_jnet_mason}

An authored two-node example isolates this point. Start from $(x_1,x_2)=(1,0)$ and define each update as copying the other node's current value. A simultaneous update from the old state gives $(0,1)$, with mean $1/2$. Sequential order $(1,2)$ gives $(0,0)$; order $(2,1)$ gives $(1,1)$. If the two sequential orders are equally likely, the mean is still $1/2$, but the variance of the resulting network mean is $1/4$, compared with zero for the deterministic simultaneous update. Agreement in an average therefore does not establish agreement in the transition law or its uncertainty. These are calculations for authored updates, not native MASON output.~\citep{exec25_jnet_choice}

Random-state identity must also be explicit. The selected seed setter narrows the supplied long integer to 32 bits, so seeds seven and $7+2^{32}$ become the same effective seed. Starting a simulation primes the generator and resets the schedule. Restoring a serialized state instead invokes a wake-from-checkpoint hook. The stream checkpoint API throws on failure, whereas convenience file wrappers print exceptions and return null. A runner should distinguish restart-from-seed from continuation-from-state and type checkpoint failure rather than interpret null as a successful empty result.~\citep{exec25_jnet_mason,exec25_jnet_choice}

The parallel-sequence documentation assigns synchronization of model data and the shared random generator to the caller. The selected RandomSequence has a switch for synchronized draws, and the fast random generator itself is unsynchronized. Locking can prevent a data race without fixing which thread consumes which random draw. Reproducible coupled experiments therefore need task-specific random streams or a deterministic allocation of draws, together with a declared update schedule and thread model. No native replay or throughput claim is made here.~\citep{exec25_jnet_mason}

MASON is a conditional simulation reference for public known-law studies. Its principal license is AFL 3.0 with named component exceptions, including the random generator's BSD terms. A proposed integration should isolate the headless simulation subset, inspect its actual distribution and test copied state, seeded starts, checkpoints, interrupted runs and independent replicates. A simulated policy result is conditional on the authored mechanism, not validation of a bank mechanism.~\citep{exec25_jnet_mason,exec25_jnet_choice}

\subsection{Repast Simphony: network direction, priority conventions and named streams}

Repast's network builder chooses directed or undirected implementations backed by ordered sparse multigraphs. Its generic adjacency method takes the union of predecessors and successors in a set, removing duplicate adjacent vertices. Separate in-edge and out-edge accessors preserve the directional distinction. Edge objects carry a mutable weight, defaulting to one. The same exposure questions therefore recur: a generic adjacent-node average cannot silently replace an incoming, weighted or transaction-multiplicity exposure.~\citep{exec25_jnet_repast}

The scheduler uses a priority convention that differs from MASON's integer ordering: larger numeric priorities execute earlier. It has explicit first, last and random-priority tokens. The inspected group delegate first establishes an order-index sequence and then intersperses random-priority actions using the random registry, or shuffles them when there are no numeric-priority actions. A port that copies numeric order values from one framework into another can reverse the intended chronology. The proposed conversion should map semantic phases, such as intervention assignment, response update and observation, before mapping numeric priorities.~\citep{exec25_jnet_repast}

The default random registry begins with a timestamp-derived integer seed. Named generators can be registered with explicit seeds. The inspected set-seed implementation invalidates default distributions, while the broader reset clears named distributions and generators and recreates the default from the clock. These distinctions are more specific than a generic statement that resetting randomness makes an experiment reproducible. A recorded seed alone also cannot establish independence of replicates or equivalence of conditional assignment laws. The runner must bind the effective generator state, name, version, draw allocation and reset path.~\citep{exec25_jnet_repast}

Repast is a conditional alternative for a richer agent-based experiment environment. The core source has BSD-style redistribution conditions and names separately licensed dependencies. The selected core manifest identifies version 2.12.0, but that does not qualify every plugin or installed runtime. The proposed choice between Repast and MASON should compare required scheduling, context, data-output and deployment features on the same public experiment, including failed requests and restart costs. A clean API response is not an identified network effect.~\citep{exec25_jnet_repast,exec25_jnet_choice}

\subsection{From simulation components to a defensible causal experiment}

To connect these components to the proposal, let $A$ denote a frozen graph, $U$ exogenous draws, $D$ an assignment and $F$ a specified state-transition system. A public simulator evaluates the consequences of $(A,U,D,F)$ under declared scheduling and observation clocks. It does not discover $F$ or establish that the observed bank assignment is independent of latent risk. Such public experiments can nevertheless test software semantics and uncertainty in regimes where the truth is known. That is the appropriate inexpensive precursor to an authentic bank study.~\citep{exec25_jnet_choice}

For two connected borrowers, independent Bernoulli assignment with probability $p$ supports all four assignments with probabilities $(1-p)^2$, $p(1-p)$, $(1-p)p$ and $p^2$. Assigning them together as one cluster supports only $(0,0)$ and $(1,1)$, with probabilities $1-p$ and $p$. At $p=1/5$, the mixed assignments have probability $4/25$ each under the independent design and zero under the single-cluster design. Neither a shared seed nor a common graph object can make those laws interchangeable. Exposure-probability and refusal tests should therefore begin with exact enumeration.~\citep{exec25_jnet_choice}

The integration plan has four independent component decisions. JGraphT supplies conditional graph operators; GraphStream supplies conditional event ingestion; MASON and Repast supply alternative public experiment environments. Direct use should preserve source contracts where they meet the scientific specification. A language port is justified only after the graph, clock, random-stream and failure contracts have an executable reference and the target runtime offers a concrete benefit. Components unsuitable for a particular target should be rejected for that use with the target-specific reason, rather than described as universally deficient.~\citep{exec25_jnet_choice}

The proposed CLI and MCP request should bind source revision and artifact, graph identity and node ordering, edge orientation and multiplicity, event and availability clocks, target population, assignment law, exposure function, update phases, effective random streams and checkpoint provenance. Its result should retain scientific assumptions, typed failures and complete costs. Small public tests can then compare directed versus incident exposure, weighted versus hop paths, induced versus background graphs, live-view versus copied state, synchronous versus sequential updates, seed collisions, restart semantics and assignment support. Native tests, full dependency rights and causal estimator qualification remain separate obligations.~\citep{exec25_jnet_choice}

Using the same hypothetical USD30,000 study budget, cost per 1,000 attempted requests is USD30, per 600 completions USD50, and per 400 accepted outputs USD75. These denominators expose the economic consequence of failure and refusal without claiming measured returns. The source review identifies concrete routes to public qualification while retaining search tails, unknown bank structure and unrestricted dependence as unresolved. It does not establish complete JVM network qualification.~\citep{exec25_jnet_choice}
